\documentclass[11pt]{article}
\usepackage[margin=1in]{geometry}
\usepackage{enumitem}
\usepackage{amsmath,amssymb}
\usepackage{hyperref}

\setlist[enumerate]{itemsep=1.1em, topsep=0.7em}

\begin{document}

\begin{center}
{\LARGE \textbf{Logic \& Reasoning — Detailed Solution Sheet}}\\[0.2em]
{\small Solutions for the 20-problem practice sheet in \texttt{Brain\_Teasers\_Reasoning\_logic.tex}}
\end{center}

\section*{Solutions}

\begin{enumerate}[label=\textbf{\arabic*.}]

% =========================================================
\item \textbf{Two jugs (water \& alcohol).}
\emph{Claim: the amount of water that ends up in the alcohol jug equals the amount of alcohol that ends up in the water jug.}

Let the alcohol jug initially contain only alcohol. Pour $Q$ water into it and mix. Now the alcohol jug contains total volume $V$ (its original volume) and within it exactly $Q$ of that is water (because we poured $Q$ water in).

Now we pour back volume $Q$ of the \emph{mixture} into the water jug.
Because the mixture is well-mixed, the fraction of water in the alcohol jug is $\frac{Q}{V}$. Therefore, the amount of water contained in the poured-back $Q$ mixture is
\[
Q\cdot \frac{Q}{V}=\frac{Q^2}{V}.
\]
So the water remaining in the alcohol jug is
\[
Q-\frac{Q^2}{V}.
\]
At the same time, the amount of \emph{alcohol} in the poured-back mixture is
\[
Q-\frac{Q^2}{V}.
\]
That alcohol is exactly what moves into the water jug.

So:
\[
\boxed{\text{Water left in the alcohol jug }=\text{ alcohol transferred to the water jug.}}
\]
Intuition: you transfer $Q$ out and $Q$ back in; whatever \emph{water} stays behind in the alcohol jug must be exactly compensated by the same volume of \emph{alcohol} that ended up in the water jug.

\hfill{\small\textit{Ref: Heard on the Street, Ch.\ 1, Q1.1}}

% =========================================================
\item \textbf{Two bells.}
If a bell rings 5 times per minute at regular intervals, its period is
\[
\frac{60\text{ sec}}{5}=12\text{ sec}.
\]
The other rings 4 times per minute, so its period is
\[
\frac{60}{4}=15\text{ sec}.
\]
They ring together again after a time equal to the least common multiple of 12 and 15:
\[
\mathrm{lcm}(12,15)=60\text{ sec}.
\]
So the next simultaneous ring occurs after \(\boxed{1\text{ minute}}\).

\hfill{\small\textit{Ref: Heard on the Street, Ch.\ 1, Q1.2}}

% =========================================================
\item \textbf{Sum of integers.}
Use the standard pairing trick:
\[
1+2+\cdots+1007=\frac{1007\cdot 1008}{2}.
\]
Compute \(1007\cdot 1008=1007(1000+8)=1{,}007{,}000+8{,}056=1{,}015{,}056\). Hence
\[
\frac{1{,}015{,}056}{2}=507{,}528.
\]
So the sum is \(\boxed{507{,}528}\).

\hfill{\small\textit{Ref: Heard on the Street, Ch.\ 1, Q1.3}}

% =========================================================
\item \textbf{Socks in the dark.}
There are only 2 colors (red/blue). In the worst case, after picking 2 socks you could have 1 red and 1 blue (no pair yet).
The \emph{third} sock must match one of these colors, guaranteeing a pair.

Thus the minimum number is \(\boxed{3}\).

\hfill{\small\textit{Ref: Heard on the Street, Ch.\ 1, Q1.15}}

% =========================================================
\item \textbf{Lily pads (doubling).}
Let the covered area (or number of pads) double each day.

Starting from 1 pad, full coverage happens on day 30. Then on day 29 it was half-covered, on day 28 quarter-covered, etc.

Starting from 8 pads means starting from \(2^3\) pads, i.e.\ you are effectively \emph{3 doubling steps ahead}.
So you reach full coverage 3 days earlier:
\[
30-3=27.
\]
Answer: \(\boxed{27\text{ days}}\).

\hfill{\small\textit{Ref: Heard on the Street, Ch.\ 1, Q1.21}}

% =========================================================
\item \textbf{Trailing zeros of \(100!\).}
A trailing zero comes from a factor of 10, i.e.\ a pair \(2\times 5\).
In \(100!\), there are many more 2s than 5s, so the number of trailing zeros equals the number of factors 5.

Count multiples of 5:
\[
\left\lfloor \frac{100}{5}\right\rfloor=20
\]
Count extra 5s coming from multiples of \(25=5^2\):
\[
\left\lfloor \frac{100}{25}\right\rfloor=4
\]
No \(5^3=125\) multiples.
Hence total 5s \(=20+4=24\), so trailing zeros \(=\boxed{24}\).

\hfill{\small\textit{Ref: Heard on the Street, Ch.\ 1, Q1.54}}

% =========================================================

\item \textbf{Macro-cube erosion.}
Total micro-cubes: \(10^3=1000\).

After removing the outer layer, what remains is the inner cube of side length \(10-2=8\), i.e.\ \(8^3=512\) micro-cubes remain attached.

Therefore, micro-cubes that fell off:
\[
1000-512=488.
\]
Answer: \(\boxed{488}\).

\hfill{\small\textit{Ref: Heard on the Street, Ch.\ 1, Q1.9}}

% =========================================================
\item \textbf{Broken clock face.}
The numbers \(1\) to \(12\) sum to
\[
1+2+\cdots+12=\frac{12\cdot 13}{2}=78.
\]
If the face breaks into 3 pieces with equal sums, each piece must sum to \(78/3=26\).

One valid partition is:
\[
\{11,12,1,2\}\quad (11+12+1+2=26),
\]
\[
\{3,4,9,10\}\quad (3+4+9+10=26),
\]
\[
\{5,6,7,8\}\quad (5+6+7+8=26).
\]
So one correct answer is:
\[
\boxed{\{11,12,1,2\},\ \{3,4,9,10\},\ \{5,6,7,8\}.}
\]
(Any partition into three groups summing to 26 is acceptable.)

\hfill{\small\textit{Ref: Heard on the Street, Ch.\ 1, Q1.4}}


% =========================================================
\item \textbf{100 bulbs / 100 brokers.}
Bulb \(k\) is toggled once for each divisor of \(k\) (because broker \(d\) toggles bulb \(k\) iff \(d\mid k\)).

A bulb ends \emph{on} iff it is toggled an odd number of times, i.e.\ iff \(k\) has an odd number of divisors.

A number has an odd number of divisors exactly when it is a perfect square (divisors pair up except the square root).

So the bulbs that remain on are:
\[
1^2,2^2,3^2,\dots,10^2
\]
i.e.\ \(1,4,9,16,25,36,49,64,81,100\). There are \(\boxed{10}\) such bulbs.

\hfill{\small\textit{Ref: Heard on the Street, Ch.\ 1, Q1.14}}

% =========================================================
\item \textbf{One fuse, 30 seconds.}
Light the fuse at \emph{both ends} simultaneously.

Even though it burns non-uniformly, the key fact is:
\begin{itemize}
\item Burning from one end takes 60 seconds total.
\item Burning from both ends means the fuse is being consumed from \emph{two sides at once}, so the time until it is fully consumed is exactly 30 seconds.
\end{itemize}
When the fuse burns out, \(\boxed{30\text{ seconds}}\) have elapsed.

\hfill{\small\textit{Ref: Heard on the Street, Ch.\ 1, Q1.51}}

% =========================================================
\item \textbf{Two fuses, 45 seconds.}
We want 45 seconds using two 60-second non-uniform fuses.

\begin{enumerate}[leftmargin=2em]
\item At time 0, light Fuse A at \emph{both ends} and Fuse B at \emph{one end}.
\item Fuse A burns out in 30 seconds (same reasoning as Problem 10).
\item Exactly then (at 30 seconds), light the \emph{other end} of Fuse B as well.
\end{enumerate}

At the moment we reach 30 seconds, Fuse B has been burning for 30 seconds from one end, so it has \emph{exactly 30 seconds of burn-time remaining} (regardless of how much physical length is left).

Lighting its other end makes the remaining burn-time halve to 15 seconds.

Total elapsed time \(=30+15=45\) seconds.

\[
\boxed{45\text{ seconds}}
\]

\hfill{\small\textit{Ref: Heard on the Street, Ch.\ 1, Q1.52}}

% =========================================================
\item \textbf{River crossing (torch).}
Times: \(D=1,\,C=2,\,B=5,\,A=10\).

A standard optimal strategy is to use the two fastest (1 and 2) as ``shuttles'' to bring the torch back.

\medskip
\noindent\textbf{Strategy (optimal):}
\begin{enumerate}[leftmargin=2em]
\item \(D\) and \(C\) cross: time \(2\). (Torch on far side.)
\item \(D\) returns: time \(+1\). (Total \(3\).)
\item \(A\) and \(B\) cross together: time \(+10\). (Total \(13\).)
\item \(C\) returns: time \(+2\). (Total \(15\).)
\item \(D\) and \(C\) cross again: time \(+2\). (Total \(17\).)
\end{enumerate}

So the minimum total time is \(\boxed{17\text{ minutes}}\).

\medskip
\noindent\textbf{Why is it optimal?}  
The slowest two \(A\) and \(B\) must cross (cost at least 10 once). To avoid having \(A\) (or \(B\)) return with the torch (very expensive), we use the fastest returns. The above plan is the classic minimax schedule and beats alternatives like sending \(A\) with \(D\) and back, which forces a slow return.

\hfill{\small\textit{Ref: A Practical Guide, Ch.\ 2, \S2.2 ``River crossing''}}

% =========================================================
\item \textbf{Coin piles (blindfolded).}
Let \(H\) be the (unknown) number of heads among the 1000 coins. We know \(H=20\), but we cannot identify them by touch.

\medskip
\noindent\textbf{Construction:}
\begin{enumerate}[leftmargin=2em]
\item Make Pile 1 by taking \emph{any} 20 coins (without looking).
\item Put the remaining 980 coins into Pile 2.
\item Flip \emph{every} coin in Pile 1.
\end{enumerate}

\medskip
\noindent\textbf{Why it works (invariant argument):}
Suppose Pile 1 initially contains \(h\) heads (and \(20-h\) tails). Then Pile 2 contains \(20-h\) heads (because there are 20 heads total).

After flipping all coins in Pile 1:
\begin{itemize}
\item those \(h\) heads become tails,
\item those \(20-h\) tails become heads,
\end{itemize}
so Pile 1 ends with \(20-h\) heads, which equals the number of heads in Pile 2.

Therefore the two piles have the same number of heads.

\[
\boxed{\text{Take 20 coins, flip them all; split is guaranteed.}}
\]

\hfill{\small\textit{Ref: A Practical Guide, Ch.\ 2, \S2.4 ``Coin piles''}}

% =========================================================
\item \textbf{Last ball (invariant).}
Let \(R\) be the number of red balls. Track the parity (even/odd) of \(R\).

Rules:
\begin{itemize}
\item If draw two blues: remove 0 red, add 1 blue $\Rightarrow R$ unchanged.
\item If draw one red and one blue: remove 1 red, add 1 red $\Rightarrow R$ unchanged.
\item If draw two reds: remove 2 red, add 1 blue $\Rightarrow R\mapsto R-2$.
\end{itemize}

So \(R\) changes only by \(-2\) in the last case. Therefore \(\boxed{R \bmod 2 \text{ is invariant}}\).

At the end there is exactly one ball:
\begin{itemize}
\item If the last ball is red, then final \(R=1\) (odd).
\item If the last ball is blue, then final \(R=0\) (even).
\end{itemize}

\medskip
\noindent\textbf{Case 1: start with 14 red (even).}  
Parity stays even, so final \(R=0\). Last ball is \(\boxed{\text{blue}}\).

\medskip
\noindent\textbf{Case 2: start with 13 red (odd).}  
Parity stays odd, so final \(R=1\). Last ball is \(\boxed{\text{red}}\).

\hfill{\small\textit{Ref: A Practical Guide, Ch.\ 2, ``Last ball''}}

% =========================================================
\item \textbf{90 coins, paid weighings.}
There are 90 coins; exactly one is counterfeit and may be \emph{heavier or lighter}. Each weighing on a balance has 3 outcomes: left heavy, right heavy, or balanced.

\medskip
\noindent\textbf{Step 1: Information-theoretic lower bound.}  
There are \(90\times 2 = 180\) possibilities (which coin, and whether heavy or light).
With \(n\) weighings, the maximum number of distinguishable outcomes is \(3^n\). So we need
\[
3^n \ge 180.
\]
We check: \(3^4=81<180\) and \(3^5=243\ge 180\).  
Thus \(\boxed{n\ge 5}\): at least 5 weighings are necessary (cost at least \$500 in the worst case).

\medskip
\noindent\textbf{Step 2: A 5-weighing strategy (balanced-ternary coding).}  
We now show 5 weighings are also sufficient, hence optimal.

Represent each coin \(i\) by a 5-digit code
\[
c_i=(c_{i1},c_{i2},c_{i3},c_{i4},c_{i5})\in\{-1,0,1\}^5,\quad c_i\neq (0,0,0,0,0),
\]
with the constraint that no two coins use opposite codes (i.e.\ never use both \(c\) and \(-c\) for two different coins), because that would create ambiguity between ``coin \(i\) heavy'' and ``coin \(j\) light''.

There are \(\frac{3^5-1}{2}=121\) available code-pairs up to sign, so we can assign 90 coins distinct codes.

\medskip
\noindent\textbf{How to perform weighing \(k\):}
\begin{itemize}
\item Put coin \(i\) on the \textbf{left} pan if \(c_{ik}=+1\).
\item Put coin \(i\) on the \textbf{right} pan if \(c_{ik}=-1\).
\item Leave coin \(i\) \textbf{off} the scale if \(c_{ik}=0\).
\end{itemize}
Choose the 90 codes so that each weighing has the same number of coins on left and right (this is always possible because you have lots of flexibility: 121 pairs available for only 90 coins).

\medskip
\noindent\textbf{Decoding:}
Encode the outcomes as a 5-vector \(o\in\{-1,0,1\}^5\):
\[
o_k=
\begin{cases}
+1 & \text{if left side is heavier on weighing }k,\\
-1 & \text{if right side is heavier on weighing }k,\\
0 & \text{if balanced.}
\end{cases}
\]
Then:
\begin{itemize}
\item If \(o=c_i\), coin \(i\) is \textbf{heavier}.
\item If \(o=-c_i\), coin \(i\) is \textbf{lighter}.
\end{itemize}
This works because a heavier counterfeit coin pushes the balance toward the side where it was placed, while a lighter one pushes the balance toward the opposite side.

\medskip
\noindent\textbf{Conclusion:} 5 weighings suffice and 4 cannot, so the optimal worst-case cost is
\[
\boxed{5 \text{ weighings } \Rightarrow \$500.}
\]

\hfill{\small\textit{Ref: Heard on the Street, Ch.\ 1, Q1.18}}

% =========================================================
\item \textbf{Motorcyclists and a fly.}
The motorcycles approach each other at relative speed \(20+30=50\) mph.
Time until they meet:
\[
t=\frac{25\ \text{miles}}{50\ \text{mph}}=0.5\ \text{hours}.
\]
The fly flies at 40 mph for the entire 0.5 hours, so distance:
\[
40\times 0.5=20\ \text{miles}.
\]
Answer: \(\boxed{20\text{ miles}}\).

\hfill{\small\textit{Ref: Heard on the Street, Ch.\ 1, Q1.28}}

% =========================================================
\item \textbf{25 horses, top 3 (5 lanes).}
\textbf{Goal: find the 3 fastest with minimum races, no stopwatch.}

\medskip
\noindent\textbf{Step 1 (5 races):} Split horses into 5 groups of 5 and race each group.  
Within each group you now know the order, e.g.\ group A: \(A_1>A_2>A_3>A_4>A_5\) (where \(A_1\) is fastest of group A).

\medskip
\noindent\textbf{Step 2 (6th race):} Race the 5 group winners \(A_1,B_1,C_1,D_1,E_1\).  
Suppose the finish order is
\[
A_1 > B_1 > C_1 > D_1 > E_1.
\]
Then:
\begin{itemize}
\item \(D_1\) and \(E_1\) cannot be in the overall top 3, because at least 3 horses (\(A_1,B_1,C_1\)) beat them.
\item In group C, only \(C_1\) can still be in the overall top 3 (since \(C_2\) is slower than \(C_1\), and already \(A_1,B_1\) beat \(C_1\)).
\item In group B, only \(B_1,B_2\) can still be in the overall top 3.
\item In group A, only \(A_1,A_2,A_3\) can still be in the overall top 3.
\end{itemize}

So the only candidates for positions 2 and 3 overall are:
\[
\{A_2,A_3,B_1,B_2,C_1\}.
\]
(\(A_1\) is already known to be the overall fastest.)

\medskip
\noindent\textbf{Step 3 (7th race):} Race these 5 candidates. The top two finishers from this race, together with \(A_1\), are the overall top 3.

\medskip
Thus \(\boxed{7}\) races suffice.

\medskip
\noindent\textbf{Why 6 races cannot guarantee it?}  
After 6 races you typically still have at least 5 plausible candidates for ranks 2 and 3, and without a 7th race you cannot fully order them. (This is the classic lower-bound argument in this puzzle.)

\hfill{\small\textit{Ref: A Practical Guide, Ch.\ 2, ``Horse race''}}

% =========================================================
\item \textbf{12 balls, heavier or lighter (3 weighings).}
\textbf{Key idea: balanced-ternary design with 3 weighings.}

Each weighing has 3 outcomes, so 3 weighings can distinguish up to \(3^3=27\) outcomes.
We need to identify 24 possibilities (12 balls $\times$ heavy/light), so 3 weighings is enough in principle.

\medskip
\noindent\textbf{Coding method.}
Assign each ball \(i\) a 3-digit code \(c_i\in\{-1,0,1\}^3\), no code is all zeros, and no two balls have opposite codes.

Interpretation per weighing \(k\):
\begin{itemize}
\item \(+1\): put ball on the left pan,
\item \(-1\): put ball on the right pan,
\item \(0\): leave it off the scale.
\end{itemize}
Outcome vector \(o\in\{-1,0,1\}^3\) is defined as:
\[
o_k=
\begin{cases}
+1 & \text{left heavy},\\
-1 & \text{right heavy},\\
0 & \text{balanced}.
\end{cases}
\]
Then:
\[
o=c_i \Rightarrow \text{ball }i\text{ is heavy},\qquad
o=-c_i \Rightarrow \text{ball }i\text{ is light}.
\]

\medskip
\noindent\textbf{One explicit working assignment (balanced each weighing).}

\begin{center}
\begin{tabular}{c|ccc}
Ball & Weigh 1 & Weigh 2 & Weigh 3\\
\hline
1  & R & R & R\\
2  & R & R & --\\
3  & R & R & L\\
4  & L & --& L\\
5  & L & --& --\\
6  & L & --& R\\
7  & R & L & R\\
8  & L & R & --\\
9  & --& L & L\\
10 & --& L & --\\
11 & --& L & R\\
12 & --& --& L\\
\end{tabular}
\end{center}

Here ``L'' means left pan, ``R'' right pan, and ``--'' means off the scale.

\medskip
\noindent\textbf{How to use it.}
Perform the three weighings exactly as specified above. Convert results to \(o\in\{-1,0,1\}^3\):
\begin{itemize}
\item left pan heavier $\to +1$,
\item right pan heavier $\to -1$,
\item balance $\to 0$.
\end{itemize}
Compare \(o\) to the code table:
\begin{itemize}
\item If \(o\) matches ball \(i\)'s code, ball \(i\) is heavy.
\item If \(o\) matches the negative of ball \(i\)'s code, ball \(i\) is light.
\end{itemize}

This uniquely identifies the defective ball and whether it is heavier or lighter in \(\boxed{3}\) weighings.

\hfill{\small\textit{Ref: A Practical Guide, Ch.\ 2, ``Defective ball''}}

% =========================================================
\item \textbf{Two glass balls, 100 floors (minimax).}
We want to minimize the \emph{worst-case} number of drops.

With 2 balls, once the first ball breaks, the second must be used for a linear search upward from the last safe floor.

So we choose floors so that the remaining worst-case trials are balanced.

Let the first drop be from floor \(n\). If it breaks, we may need up to \(n-1\) more drops (linear search) $\Rightarrow$ worst-case \(n\).
If it survives, next drop should be \(n+(n-1)\), so that if it breaks there, we have at most \(n-1\) remaining linear checks, giving worst-case still \(n\).

Continue with decreasing steps:
\[
n,\quad n+(n-1),\quad n+(n-1)+(n-2),\ \dots
\]
We need the sum \(n+(n-1)+\cdots+1=\frac{n(n+1)}{2}\) to be at least 100.

Solve:
\[
\frac{n(n+1)}{2}\ge 100.
\]
Check \(n=13\): \(13\cdot 14/2=91<100\).  
Check \(n=14\): \(14\cdot 15/2=105\ge 100\).

So the optimal minimax number is \(\boxed{14}\) drops.

\medskip
\noindent\textbf{Concrete schedule:}
Drop ball 1 from floors:
\[
14,\ 27,\ 39,\ 50,\ 60,\ 69,\ 77,\ 84,\ 90,\ 95,\ 99,\ 100
\]
(using step sizes 14,13,12,11,10,9,8,7,6,5,4,1).
When it breaks, use ball 2 to search linearly between the last safe floor + 1 and the breaking floor - 1.

Worst case: 14 drops.

\hfill{\small\textit{Ref: A Practical Guide, Ch.\ 2, ``Glass ball''}}

% =========================================================
\item \textbf{50 wise men and a glass.}
\textbf{Classic ``counter'' strategy (guarantees eventual freedom).}

\medskip
\noindent\textbf{Plan:} Choose one wise man as the \emph{counter} (say Wise Man 1). The other 49 are \emph{signalers}.

The glass has two states. Call them:
\[
\text{Upright} = 0,\qquad \text{Upside-down} = 1.
\]

\medskip
\noindent\textbf{Rules agreed beforehand.}

\paragraph{Counter (Wise Man 1).}
\begin{itemize}
\item On his \emph{very first} visit, he sets the glass to Upright (0) and sets his internal count to 0 (no declaration).
\item On any later visit:
  \begin{itemize}
  \item If the glass is Upside-down (1), he flips it back to Upright (0) and increments his count by 1.
  \item If the glass is Upright (0), he does nothing.
  \end{itemize}
\item When his count reaches 49, the next time he is called he declares: \emph{``All 50 wise men have been called at least once.''}
\end{itemize}

\paragraph{Each signaler (Wise Men 2--50).}
\begin{itemize}
\item Keep track of whether you have been called before.
\item If this is \emph{your first time} being called and the glass is Upright (0), then flip it to Upside-down (1).
\item Otherwise (if it is not your first time, or the glass is already Upside-down), do nothing.
\end{itemize}

\medskip
\noindent\textbf{Why this works (key invariants).}
\begin{itemize}
\item The glass being Upside-down acts like a ``pending message'' that \emph{some new signaler} has been called and has not yet been counted.
\item While the glass is Upside-down, no other signaler will flip it (by rule), so signals cannot stack up and cannot be lost.
\item The counter is the only person who resets the glass from 1 back to 0, and each reset corresponds to counting exactly one new signaler.
\end{itemize}

\medskip
\noindent\textbf{Why freedom is guaranteed (eventuality).}
Assuming the sultan keeps calling people indefinitely, every signaler is called infinitely often with probability 1. In particular, each signaler will eventually be called at a time when the glass is Upright, and will then flip it to Upside-down exactly once (the first time they succeed in doing so).

Each time that happens, sooner or later the counter will also be called; he will see the glass Upside-down, flip it back, and increment his count. Hence the counter’s count eventually reaches 49.

At that moment, he knows:
\[
\text{(49 signalers have been called at least once)}\quad + \quad \text{(I have been called too)}
\]
so all 50 have been called. Therefore his declaration is correct and everyone goes free.

\[
\boxed{\text{Strategy guarantees eventual freedom.}}
\]

\hfill{\small\textit{Ref: A Practical Guide, ``Wise men''}}

\end{enumerate}

\end{document}